% !TeX root = ../../Book5.tex
\chapter*{后记}
\addcontentsline{toc}{chapter}{后记}

本卷从有限群的线性作用出发，最后走到不变量环、代数群与微分方程。回看这些内容，表示论让我们能够换一种方式认识代数对象：同一个群可以通过特征标、张量积或坐标函数来研究；同一个代数可以通过简单模、投射模或箭图中的线性映射来研究。每换一种表示，原先隐蔽的关系便可能成为一条可以计算的公式。

\ProseParagraphBreak
这种转变也有代价。特征标看得见合成因子，却不总能看见扩张；维数向量记录各顶点的大小，却不能独自决定箭图表示；相同的根系也未必决定实际基域上的群。不变量所能区分的对象、一个函子所遗忘的数据、一个公式成立所需的条件，都是本卷反复要求读者检查的事。希望这些提醒最终成为习惯，而不只是定理旁的一行限制。

\ProseParagraphBreak
几项较长的证明值得在读完后再回顾。对称群的标准基把组合重写和线性无关联系起来；箭图的反射把维数变换变成表示之间的操作；最高权分类先构造模，再证明有限维，特征标公式才随后计算它。它们都不只是得出结论的方法，也解释了为什么那些参数、图形和整数会出现在答案中。把一个证明重新用于最小的非平凡例子，往往比继续记住更多定理名称更有收获。

\ProseParagraphBreak
本书没有试图把表示论的每个方向都压缩成一章。模表示的深层块论、一般代数群的算术形式、几何表示论、量子群与范畴化，以及解析和微分 Galois 理论，都有超出本卷的内容。附录选取定义、代表性结果和可以核验的模型，较深的结果给出注明输入与省略步骤的证明概要，并提供准确出处。读者若准备深入某一方向，应沿相应资料学习完整理论，也应重新检查那一领域对底域、拓扑、光滑性及有限性的约定。

\ProseParagraphBreak
五卷的写作到这里告一段落，修订仍将继续。一个证明有没有说清当前使用的是哪条结论，一个定义能否独立查阅，一道练习是否真正检验了某个新的判断，这些问题都需要在教学、阅读与反复计算中继续回答。愿本书能让读者带走一些可靠的方法，也留下值得自己追问的问题。
